\chapter*{Abstract}
\addcontentsline{toc}{chapter}{Abstract}
\markboth{Abstract}{Abstract}

In LLM-based program repair, each round has two costs: the model call that
proposes a patch and the test run that checks it. Known techniques make the
test run cheap: the patched file can be compiled in process, the test runtime
kept alive, and the tests that cannot reach the change skipped. What remains is
the model call. A strong model repairs more bugs but costs about twenty times
more per token than a cheap one, and a repair bot must pick a model before it
knows how hard the bug is. Starting at the strong model pays its price for easy
bugs. Staying at the cheap model burns its budget on bugs it cannot repair. And
at a switch between models, what the first model learned is either dropped or
carried in a form that may bias the second.

This proposal builds a repair bot that tries the cheap model first and
\emph{steps up} to the strong model only when the tests show that the session
has stalled. A \emph{Warm Substrate} pays the build and the start-up of the
test runtime once per session instead of once per candidate, so the cheap
model's extra candidates cost little to check. An \emph{Evidence Memory} keeps
every verdict, failing test and failure signature, and screens the candidates
of both models. An \emph{Escalation Controller} applies a rule, \Step, that
moves the session up a tier after a verified stall and writes a \emph{handover
note} from the evidence. A \emph{Budget Guard} stops sessions that are stuck or
out of budget and reports why. With a cold handover and an unchanged budget for
the strong model, the bot repairs every bug that the strong model alone repairs
on the same recorded sessions; what it changes is the cost.

The evaluation uses $100$ Defects4J bugs. Sessions of each model alone are
recorded once. Every escalation policy is then replayed on the same candidates
and compared with a router that knows the outcome in hindsight. Live runs then
measure what the handover note is worth to the strong model and, on bugs held
out from that choice, whether the whole system lowers the cost per repaired
bug. The saving is expected mainly in money and tokens; in time it depends on
how long the cheap tier delays the bugs it cannot repair.

\vspace{1em}
\noindent\textbf{Keywords:} automated program repair, large language models,
model escalation, patch validation, cost analysis
